\chapter{Events and the Command Dispatcher}
\label{ch:events-dispatcher}

In Chapter~\ref{ch:attributes-props} we learned how props become DOM configuration.
Now we hit the other half of interactivity: \textbf{input flowing back into the system}.

Relax uses the word ``event'' in two different senses, and it helps to separate them in your head:

\begin{itemize}
  \item \textbf{DOM events}: \texttt{click}, \texttt{input}, \texttt{change}, \dots created by the browser and delivered to listeners on real elements.
  \item \textbf{Commands}: named messages with payloads delivered through an in-memory bus.
\end{itemize}

These are implemented as two tiny modules:

\begin{itemize}
  \item \texttt{src/events.ts} --- converts Relax's \texttt{props.on} map into real DOM listeners with predictable binding.
  \item \texttt{src/dispatcher.ts} --- a deterministic pub/sub for named commands.
\end{itemize}

And, as usual in Relax, the real contract is pinned down by tests:

\begin{itemize}
  \item \texttt{src/\_\_tests\_\_/events.test.ts}
  \item \texttt{src/\_\_tests\_\_/dispatcher.test.ts}
\end{itemize}

\section{DOM events: the two problems every renderer must solve}

If a UI library wants to support authoring like:

\begin{lstlisting}[language=JavaScript]
<button onClick={this.onClick}>Hit</button>
\end{lstlisting}

it has to solve two very unglamorous problems:

\begin{enumerate}
  \item \textbf{Binding:} when the browser calls the handler, what should \texttt{this} be?
  \item \textbf{Removal by identity:} when the tree unmounts (or events update), can we remove exactly what we added?
\end{enumerate}

The second one is the bigger deal in practice.
The DOM removes listeners by \emph{function identity}, not by comparing source code or ``equivalent'' functions.

Relax's answer is: \textbf{always register a wrapper listener}. The wrapper is what gets stored and later removed.

\subsection*{Diagram: why wrapper identity matters for removal}
\begin{center}
\begin{tikzpicture}[
  node distance=10mm and 14mm,
  box/.style={draw, rounded corners, align=left, inner sep=6pt, fill=RelaxLight},
  arrow/.style={-Latex, thick}
]
\node[box] (user) {User handler\\\texttt{handler(event)}};
\node[box, right=of user] (wrap) {Relax wrapper\\\texttt{listener = (e) => handler.call(host?, e)}};
\node[box, right=of wrap] (dom) {DOM element\\\texttt{el.addEventListener(name, listener)}};

\node[box, below=of dom] (remove) {Cleanup\\\texttt{el.removeEventListener(name, listener)}};

\draw[arrow] (user) -- node[midway, above] {not registered} (wrap);
\draw[arrow] (wrap) -- node[midway, above] {registered by identity} (dom);
\draw[arrow] (dom) -- (remove);

\draw[arrow] (user) to[bend right=25] node[midway, below] {can\textbf{not} remove} (remove);
\end{tikzpicture}
\end{center}

The DOM can only remove listeners by the exact function object identity that was registered. That's why \texttt{events.ts} returns the wrapper.

\section{DOM event listener wiring (\texttt{src/events.ts})}
The central helper is:

\begin{quote}
	exttt{addEventListener(eventName, handler, el, hostComponent?)}
\end{quote}

\subsection{Contract}
\begin{itemize}
  \item \textbf{Inputs:}
    \begin{itemize}
      \item \texttt{eventName}: a DOM event name like \texttt{"click"}.
      \item \texttt{handler}: a function that should run when the event fires.
      \item \texttt{el}: the real DOM element.
      \item \texttt{hostComponent}: optional object used as the \texttt{this} binding.
    \end{itemize}
  \item \textbf{Output:}
    \begin{itemize}
      \item returns the concrete \texttt{listener} function object that was actually registered with \texttt{el.addEventListener}.
    \end{itemize}
\end{itemize}

Returning the registered listener is not a convenience; it's required for correctness.

\subsection{Implementation: the wrapper listener (real code)}
This is the actual implementation from \texttt{src/events.ts}. Notice the wrapper chooses between \texttt{handler.call(hostComponent, event)} and \texttt{handler(event)}:

\begin{lstlisting}[language=JavaScript]
export function addEventListener(
  eventName: string,
  handler: (payload: unknown) => void,
  el: Element,
  hostComponent: any = null
) {
  const listener = (event: Event) => {
    if (hostComponent) {
      handler.call(hostComponent, event)
    } else {
      handler(event)
    }
  }

  el.addEventListener(eventName, listener)
  return listener
}
\end{lstlisting}

This is one of those renderer tricks that's simple but non-negotiable.
If you ever ``refactor'' this into returning \texttt{handler} directly, unmount will leak listeners.

This immediately gives Relax three desirable properties:

\begin{itemize}
  \item \textbf{Stable binding semantics:} the wrapper decides whether to call with \texttt{handler.call(hostComponent, event)}.
  \item \textbf{No mutation of user handlers:} the original function object is untouched.
  \item \textbf{Reliable cleanup:} the wrapper function object can be stored and later removed.
\end{itemize}

\subsection*{Diagram: event wiring lifecycle across mount, patch, destroy}
\begin{center}
\begin{tikzpicture}[
  node distance=10mm and 12mm,
  box/.style={draw, rounded corners, align=left, inner sep=6pt, fill=RelaxLight},
  arrow/.style={-Latex, thick}
]
\node[box] (props) {VNode props\\\texttt{props.on = \{ click: fn, ... \}}};
\node[box, right=of props] (mount) {Mount\\\texttt{addEventListeners(...)}\\returns wrappers map};
\node[box, right=of mount] (store) {Store on vnode\\\texttt{vdom.listeners = wrappers}};
\node[box, below=of store] (patch) {Patch\\compare old/new\\update listeners (swap wrappers)};
\node[box, left=of patch] (destroy) {Destroy\\\texttt{removeEventListeners(vdom.listeners, el)}};

\draw[arrow] (props) -- (mount);
\draw[arrow] (mount) -- (store);
\draw[arrow] (store) -- (patch);
\draw[arrow] (store) -- (destroy);
\end{tikzpicture}
\end{center}

Even if patch-time logic changes, notice the invariant: destroy always removes \emph{whatever wrappers are currently stored on the vnode}.

\subsection{Multiple listeners: \texttt{addEventListeners(events, el, hostComponent?)}}
Mounting and patching usually deal with an event map (a record of event name to handler). Relax provides a convenience:

\begin{quote}
    exttt{addEventListeners(events, el, hostComponent)}
\end{quote}

Implementation is straightforward:

\begin{itemize}
  \item iterate \texttt{Object.entries(events)}
  \item call \texttt{addEventListener} for each
  \item store returned wrappers in \texttt{listeners: Record<string, EventListener>}
\end{itemize}

The important artifact is the returned \texttt{listeners} map. In \texttt{mount-dom.ts}, Relax stores it on the vnode (\texttt{vdom.listeners}) so that unmounting can remove exactly what was added.

\subsection{Cleanup symmetry: \texttt{removeEventListeners(listeners, el)}}
Removal is the mirror image:

\begin{lstlisting}[language=JavaScript]
Object.entries(listeners).forEach(([eventName, listener]) => {
  el.removeEventListener(eventName, listener)
})
\end{lstlisting}

Note what's \emph{not} present: the original user handlers. Only wrappers can be removed.

\section{Tests for DOM event semantics (\texttt{src/\_\_tests\_\_/events.test.ts})}
The tests are small, but they lock down the two things that tend to regress when you wrap handlers.

\subsection{Invariant 1: wrappers don't break sync or async handlers}
	exttt{events.test.ts} contains:

\begin{itemize}
  \item \texttt{synchronous event without host component, called with argument}
  \item \texttt{asynchronous event without host component}
\end{itemize}

The async test uses \texttt{flushPromises()} (\texttt{src/utils/promises.ts}) after \texttt{btn.click()} to prove that:

\begin{itemize}
  \item the wrapper does not swallow async behavior
  \item the handler is actually invoked and can \texttt{await} microtasks
\end{itemize}

This matters because event wrappers are a common source of subtle bugs where promises aren't awaited or exceptions get lost.

\subsection{Invariant 2: with a host component, \texttt{this} is correct}
Two tests cover binding (sync and async):

\begin{itemize}
  \item \texttt{synchronous event with host component, the handler is bound to component}
  \item \texttt{asynchronous event with host component, the handler is bound to component}
\end{itemize}

Both assert:

\begin{itemize}
  \item \texttt{actualBinding === hostComponent}
  \item the handler argument is an \texttt{Event}
\end{itemize}

If you remember only one idea from this chapter, make it this: \textbf{binding and cleanup are owned by the wrapper function.}

\section{The command dispatcher (\texttt{src/dispatcher.ts})}
Where \texttt{events.ts} is about \emph{browser integration}, \texttt{Dispatcher} is an internal message bus.
I use it for things like devtools hooks, app-level commands, and decoupled communication that doesn't want to bounce through the DOM.

\subsection*{Diagram: dispatcher data structures and dispatch flow}
\begin{center}
\begin{tikzpicture}[
  node distance=10mm and 12mm,
  box/.style={draw, rounded corners, align=left, inner sep=6pt, fill=RelaxLight},
  arrow/.style={-Latex, thick}
]
\node[box] (subs) {\texttt{\#subs: Map<string, Handler[]>}};
\node[box, right=of subs] (after) {\texttt{\#afterHandlers: (() => void)[]}};
\node[box, below=of subs] (dispatch) {\texttt{dispatch(name, payload)}};
\node[box, below=of dispatch] (handlers) {Run handlers for \texttt{name}\\or warn if none};
\node[box, right=of handlers] (afterRun) {Always run\\after-handlers};

\draw[arrow] (subs) -- (dispatch);
\draw[arrow] (dispatch) -- (handlers);
\draw[arrow] (handlers) -- (afterRun);
\draw[arrow] (after) -- (afterRun);
\end{tikzpicture}
\end{center}

\subsection{A minimal contract}
	exttt{Dispatcher} has three public methods:

\begin{itemize}
  \item \texttt{subscribe(commandName, handler)} \textrightarrow{} returns \texttt{unsubscribe()}
  \item \texttt{afterEveryCommand(handler)} \textrightarrow{} returns \texttt{unsubscribe()}
  \item \texttt{dispatch(commandName, payload)}
\end{itemize}

Internally it stores:

\begin{itemize}
  \item \texttt{\#subs}: \texttt{Map<string, Array<(payload: unknown) => void>>}
  \item \texttt{\#afterHandlers}: \texttt{Array<() => void>}
\end{itemize}

So you can subscribe to a particular command name (payload-aware) and/or attach a hook that runs after every dispatch (payload-agnostic).

\subsection{Subscription semantics: dedupe by function identity}
An important policy choice appears in \texttt{subscribe()}:

\begin{quote}
  registering the same handler function twice is a no-op.
\end{quote}

In code, that's \texttt{handlers.includes(handler)}.

This makes the dispatcher deterministic and prevents accidental handler explosions if a component subscribes in a loop.

\subsection{Dispatch semantics: warn on unhandled commands, always run after-hooks}
	exttt{dispatch(commandName, payload)} does:

\begin{enumerate}
  \item if there are handlers for \texttt{commandName}, call them with \texttt{payload}
  \item otherwise, \texttt{console.warn} (if available): \texttt{No handlers for command: <name>}
  \item always run \texttt{\#afterHandlers}
\end{enumerate}

That last rule is important: after-hooks run even if there were no command handlers.
It allows patterns like global logging, tracing, or devtools hooks.

\section{Tests for dispatcher semantics (\texttt{src/\_\_tests\_\_/dispatcher.test.ts})}
The dispatcher tests are a clean, readable spec for the API:

\begin{itemize}
  \item \textbf{subscribe/unsubscribe works:} the handler is called once, then not called after unsubscribe.
  \item \textbf{dedupe works:} subscribing the same handler multiple times still results in one call per dispatch.
  \item \textbf{multiple handlers work:} two different handlers both fire, and both can unsubscribe.
  \item \textbf{afterEveryCommand works:} runs once per dispatch and unsubscribes cleanly.
\end{itemize}

Notice how these tests avoid implementation details (maps, arrays, indices) and instead assert the observable contract.

\section{Common pitfalls and edges}
\begin{itemize}
  \item \textbf{Listener removal requires wrapper identity.} Never attempt to remove a handler by passing the original user function if you registered a wrapper.
  \item \textbf{Dedupe is per-function identity.} If you pass a new \texttt{() => ...} lambda each time, the dispatcher will treat it as a different handler.
  \item \textbf{Dispatcher is synchronous.} If a handler throws, that throw will propagate from \texttt{dispatch()}.
\end{itemize}

\section{Mini-exercises}
\begin{enumerate}
  \item In \texttt{src/events.ts}, add a new helper \texttt{updateEventListeners(old, next, el, hostComponent)} that returns the new wrappers map without leaking old listeners. (Don't implement it here; design the contract and write the tests first.)
  \item Extend \texttt{dispatcher.test.ts} with a test that proves after-handlers run even when there are \emph{no} command handlers for a dispatch.
  \item Thought exercise: should \texttt{Dispatcher.subscribe()} remove the command entry from \texttt{\#subs} when the last handler unsubscribes? What trade-offs does that create?
\end{enumerate}

\section{Where this plugs into the renderer}
You've already seen the integration point in Chapter~\ref{ch:mounting}:

\begin{itemize}
  \item element vnodes store \texttt{props.on} (after TSX mapping)
  \item mounting calls \texttt{addEventListeners(...)} and saves wrapper functions on \texttt{vdom.listeners}
  \item destroying later calls \texttt{removeEventListeners(vdom.listeners, el)}
\end{itemize}

The remaining questions, how listeners are updated during patch without leaking, and how host binding is preserved across updates, live in \texttt{src/patch-dom.ts} and \texttt{src/destroy-dom.ts}.

\section{Online resources}
\begin{itemize}
  \item DOM events and bubbling (MDN): \href{https://developer.mozilla.org/en-US/docs/Learn_web_development/Core/Scripting/Event_bubbling}{\url{https://developer.mozilla.org/en-US/docs/Learn_web_development/Core/Scripting/Event_bubbling}}
  \item \texttt{EventTarget.addEventListener} / \texttt{removeEventListener} (MDN): \href{https://developer.mozilla.org/en-US/docs/Web/API/EventTarget/addEventListener}{\url{https://developer.mozilla.org/en-US/docs/Web/API/EventTarget/addEventListener}}
  \item Function \texttt{call} and \texttt{this} binding (MDN): \href{https://developer.mozilla.org/en-US/docs/Web/JavaScript/Reference/Global_Objects/Function/call}{\url{https://developer.mozilla.org/en-US/docs/Web/JavaScript/Reference/Global_Objects/Function/call}}
  \item Message bus / pub-sub as a pattern (Wikipedia overview): \href{https://en.wikipedia.org/wiki/Publish%E2%80%93subscribe_pattern}{\url{https://en.wikipedia.org/wiki/Publish-subscribe_pattern}}
\end{itemize}
